% sections/performance.tex
% Included from main.tex via % sections/performance.tex
% Included from main.tex via % sections/performance.tex
% Included from main.tex via % sections/performance.tex
% Included from main.tex via \input{sections/performance}

\section{Performance Evaluation}
\label{sec:perf}

This section evaluates CARLA-Air under representative joint air-ground
workloads across three experiments: frame-rate and resource scaling
(Section~\ref{sec:perf:fps}), memory stability under sustained operation
(Section~\ref{sec:perf:vram}), and communication latency
(Section~\ref{sec:perf:rpc}).
Full configuration parameters and raw data are deferred to
Appendix~\ref{app:config}.

\paragraph{Reference hardware.}
All measurements are collected on a workstation equipped with an NVIDIA
RTX~A4000 (16\,GB GDDR6), AMD Ryzen~7 5800X (8-core, 4.7\,GHz), and 32\,GB
DDR4-3200, running Ubuntu~20.04~LTS.
The simulator runs in Epic quality mode with Town10HD loaded unless stated
otherwise.
All aerial experiments use the built-in SimpleFlight controller with default
PID gains.
CPU affinity and GPU power limits are left at system defaults to reflect
realistic research deployment conditions.

%% ---------------------------------------------------------------
\subsection{Benchmark Methodology}
\label{sec:perf:method}

Reliable performance measurement requires eliminating startup
transients---map-loading jitter, first-frame shader compilation, and actor
lifecycle initialization must be discarded before steady-state sampling begins.
Algorithm~\ref{alg:bench} formalizes the benchmark harness used throughout this
section.

\input{figures/alg_bench}

Each profile uses $T_w = 200$ warm-up ticks and $T_m = 2{,}000$ measurement
ticks.
VRAM is sampled every 60\,s.
All reported frame rates are the harmonic mean of $\mathbf{f}$---the
appropriate central tendency for rate quantities---with standard deviations
alongside.
The latency benchmark issues 500 warm-up calls followed by 5\,000 measurement
calls; actor spawn calls are each paired with an immediate destroy to prevent
scene-state accumulation from contaminating subsequent measurements.

%% ---------------------------------------------------------------
\subsection{Experiment~1: Frame Rate and Resource Scaling}
\label{sec:perf:fps}

\paragraph{Workload design.}
Under synchronous-mode operation, per-tick wall time is bounded by the slowest
of three concurrent contributors: the rendering thread (GPU-bound), the
ground-actor dispatch loop (CPU-bound), and the aerial physics engine
(CPU-bound, asynchronous at ${\approx}1{,}000$\,Hz on a dedicated thread).
Sensor rendering dominates at high resolution due to GPU memory bandwidth
saturation; actor population contributes via per-mesh draw calls.
These relationships motivate the workload stratification in
Table~\ref{tab:fps_scaling}.

\input{figures/tab_fps_scaling}

\paragraph{Analysis.}
The moderate joint configuration---combining ground traffic, an aerial agent,
and a full sensor suite---sustains $19.8 \pm 1.1$\,FPS, sufficient for
closed-loop policy evaluation at standard RL episode lengths.
Comparing the standalone ground baseline ($28.4$\,FPS) against the moderate
joint profile ($19.8$\,FPS) quantifies total integration overhead at
$8.6$\,FPS (30.3\%), of which $2.1$\,FPS is attributable to ground co-hosting
and the remaining $6.5$\,FPS to the aerial physics engine.
Crucially, this aerial overhead manifests entirely in CPU utilization ($54\%$
vs.\ $38\%$) rather than GPU memory: VRAM differs by only 39\,MiB between
ground-only and moderate-joint profiles.
The traffic surveillance profile retains comparable throughput ($20.1$\,FPS)
despite doubling the vehicle count, confirming that sensor rendering---not
actor population---is the dominant cost driver at $1920\times1080$.
The 3-hour endurance profile ($19.7 \pm 1.3$\,FPS) confirms that throughput
does not degrade under sustained operation.

%% ---------------------------------------------------------------
\subsection{Experiment~2: Memory Stability}
\label{sec:perf:vram}

\paragraph{Steady-state VRAM profile.}
The base map asset load accounts for approximately 3\,702\,MiB at idle
(Town10HD, Epic quality); each additional ground sensor contributes
4--10\,MiB at $1280\times720$.
A transient peak of approximately 5\,000\,MiB occurs during map loading and
resolves within 30\,s.
All steady-state profiles remain within 3\,878\,MiB, retaining approximately
12\,506\,MiB (76\%) of the 16\,GB device budget for co-located workloads such
as GPU-based policy training.

\paragraph{Leak verification.}
Table~\ref{tab:stability} summarises the 3-hour endurance run across 357 actor
spawn-and-destroy cycles.
A linear regression of VRAM against cycle index yields a slope of
$0.49$\,MiB/cycle with $R^2 = 0.11$, indicating no statistically significant
accumulation trend.
Fig.~\ref{fig:vram_timeseries} visualises the VRAM trace: the negligible
early-to-late drift ($3{,}868 \to 3{,}878$\,MiB, approximately 10\,MiB over
327 cycles) is attributable to residual render-target caching rather than
lifecycle leakage.

\input{figures/fig_vram_timeseries}
\input{figures/tab_stability}

Zero API errors and zero simulation crashes across 357 cycles validate the
robustness of the joint environment under the repeated agent-reset patterns
typical of reinforcement-learning training.

%% ---------------------------------------------------------------
\subsection{Experiment~3: Communication Latency}
\label{sec:perf:rpc}

Both simulation APIs operate within the same process on the loopback interface,
eliminating inter-process serialization overhead.
The two RPC servers use distinct wire protocols and port assignments (detailed
in Appendix~\ref{app:config}).
Table~\ref{tab:rpc_latency} reports round-trip latency for representative API
calls.

\input{figures/tab_rpc_latency}

\paragraph{Analysis.}
Lightweight state queries complete in 280--490\,$\mu$s, well below the per-tick
budget at 20\,FPS (50\,ms), confirming that API overhead does not contribute
meaningfully to frame-time variance.
Actor spawn latency (1\,850\,$\mu$s) reflects a one-time GPU synchronization
point for render-asset registration at episode reset.
Image capture latency (3\,200\,$\mu$s) covers the full sensor
pipeline---rendering, buffer readback, and serialization---and is overlapped
with the rendering thread at synchronous tick rates.
All measured values fall below the lower bound of bridge-based cross-process
synchronization costs reported in~\cite{transimhub}
(1\,000--5\,000\,$\mu$s per frame).

\paragraph{Tick-rate reconciliation.}
The aerial physics engine advances at ${\approx}1{,}000$\,Hz on its dedicated
thread, while the rendering tick runs at ${\approx}20$\,Hz under the moderate
joint workload.
Sensor callbacks read the aerial state at rendering tick boundaries, so each
ground-aerial sensor frame pair reflects the drone's integrated physical state
over ${\approx}50$ aerial physics steps.
Ground actor states are also resolved at tick boundaries, ensuring temporal
co-registration across all sensor modalities within a single tick.
Applications requiring finer aerial state resolution should reduce the fixed
simulation delta-time accordingly; the throughput trade-off follows from
Table~\ref{tab:fps_scaling}.

\section{Performance Evaluation}
\label{sec:perf}

This section evaluates CARLA-Air under representative joint air-ground
workloads across three experiments: frame-rate and resource scaling
(Section~\ref{sec:perf:fps}), memory stability under sustained operation
(Section~\ref{sec:perf:vram}), and communication latency
(Section~\ref{sec:perf:rpc}).
Full configuration parameters and raw data are deferred to
Appendix~\ref{app:config}.

\paragraph{Reference hardware.}
All measurements are collected on a workstation equipped with an NVIDIA
RTX~A4000 (16\,GB GDDR6), AMD Ryzen~7 5800X (8-core, 4.7\,GHz), and 32\,GB
DDR4-3200, running Ubuntu~20.04~LTS.
The simulator runs in Epic quality mode with Town10HD loaded unless stated
otherwise.
All aerial experiments use the built-in SimpleFlight controller with default
PID gains.
CPU affinity and GPU power limits are left at system defaults to reflect
realistic research deployment conditions.

%% ---------------------------------------------------------------
\subsection{Benchmark Methodology}
\label{sec:perf:method}

Reliable performance measurement requires eliminating startup
transients---map-loading jitter, first-frame shader compilation, and actor
lifecycle initialization must be discarded before steady-state sampling begins.
Algorithm~\ref{alg:bench} formalizes the benchmark harness used throughout this
section.

\input{figures/alg_bench}

Each profile uses $T_w = 200$ warm-up ticks and $T_m = 2{,}000$ measurement
ticks.
VRAM is sampled every 60\,s.
All reported frame rates are the harmonic mean of $\mathbf{f}$---the
appropriate central tendency for rate quantities---with standard deviations
alongside.
The latency benchmark issues 500 warm-up calls followed by 5\,000 measurement
calls; actor spawn calls are each paired with an immediate destroy to prevent
scene-state accumulation from contaminating subsequent measurements.

%% ---------------------------------------------------------------
\subsection{Experiment~1: Frame Rate and Resource Scaling}
\label{sec:perf:fps}

\paragraph{Workload design.}
Under synchronous-mode operation, per-tick wall time is bounded by the slowest
of three concurrent contributors: the rendering thread (GPU-bound), the
ground-actor dispatch loop (CPU-bound), and the aerial physics engine
(CPU-bound, asynchronous at ${\approx}1{,}000$\,Hz on a dedicated thread).
Sensor rendering dominates at high resolution due to GPU memory bandwidth
saturation; actor population contributes via per-mesh draw calls.
These relationships motivate the workload stratification in
Table~\ref{tab:fps_scaling}.

\input{figures/tab_fps_scaling}

\paragraph{Analysis.}
The moderate joint configuration---combining ground traffic, an aerial agent,
and a full sensor suite---sustains $19.8 \pm 1.1$\,FPS, sufficient for
closed-loop policy evaluation at standard RL episode lengths.
Comparing the standalone ground baseline ($28.4$\,FPS) against the moderate
joint profile ($19.8$\,FPS) quantifies total integration overhead at
$8.6$\,FPS (30.3\%), of which $2.1$\,FPS is attributable to ground co-hosting
and the remaining $6.5$\,FPS to the aerial physics engine.
Crucially, this aerial overhead manifests entirely in CPU utilization ($54\%$
vs.\ $38\%$) rather than GPU memory: VRAM differs by only 39\,MiB between
ground-only and moderate-joint profiles.
The traffic surveillance profile retains comparable throughput ($20.1$\,FPS)
despite doubling the vehicle count, confirming that sensor rendering---not
actor population---is the dominant cost driver at $1920\times1080$.
The 3-hour endurance profile ($19.7 \pm 1.3$\,FPS) confirms that throughput
does not degrade under sustained operation.

%% ---------------------------------------------------------------
\subsection{Experiment~2: Memory Stability}
\label{sec:perf:vram}

\paragraph{Steady-state VRAM profile.}
The base map asset load accounts for approximately 3\,702\,MiB at idle
(Town10HD, Epic quality); each additional ground sensor contributes
4--10\,MiB at $1280\times720$.
A transient peak of approximately 5\,000\,MiB occurs during map loading and
resolves within 30\,s.
All steady-state profiles remain within 3\,878\,MiB, retaining approximately
12\,506\,MiB (76\%) of the 16\,GB device budget for co-located workloads such
as GPU-based policy training.

\paragraph{Leak verification.}
Table~\ref{tab:stability} summarises the 3-hour endurance run across 357 actor
spawn-and-destroy cycles.
A linear regression of VRAM against cycle index yields a slope of
$0.49$\,MiB/cycle with $R^2 = 0.11$, indicating no statistically significant
accumulation trend.
Fig.~\ref{fig:vram_timeseries} visualises the VRAM trace: the negligible
early-to-late drift ($3{,}868 \to 3{,}878$\,MiB, approximately 10\,MiB over
327 cycles) is attributable to residual render-target caching rather than
lifecycle leakage.

\input{figures/fig_vram_timeseries}
\input{figures/tab_stability}

Zero API errors and zero simulation crashes across 357 cycles validate the
robustness of the joint environment under the repeated agent-reset patterns
typical of reinforcement-learning training.

%% ---------------------------------------------------------------
\subsection{Experiment~3: Communication Latency}
\label{sec:perf:rpc}

Both simulation APIs operate within the same process on the loopback interface,
eliminating inter-process serialization overhead.
The two RPC servers use distinct wire protocols and port assignments (detailed
in Appendix~\ref{app:config}).
Table~\ref{tab:rpc_latency} reports round-trip latency for representative API
calls.

\input{figures/tab_rpc_latency}

\paragraph{Analysis.}
Lightweight state queries complete in 280--490\,$\mu$s, well below the per-tick
budget at 20\,FPS (50\,ms), confirming that API overhead does not contribute
meaningfully to frame-time variance.
Actor spawn latency (1\,850\,$\mu$s) reflects a one-time GPU synchronization
point for render-asset registration at episode reset.
Image capture latency (3\,200\,$\mu$s) covers the full sensor
pipeline---rendering, buffer readback, and serialization---and is overlapped
with the rendering thread at synchronous tick rates.
All measured values fall below the lower bound of bridge-based cross-process
synchronization costs reported in~\cite{transimhub}
(1\,000--5\,000\,$\mu$s per frame).

\paragraph{Tick-rate reconciliation.}
The aerial physics engine advances at ${\approx}1{,}000$\,Hz on its dedicated
thread, while the rendering tick runs at ${\approx}20$\,Hz under the moderate
joint workload.
Sensor callbacks read the aerial state at rendering tick boundaries, so each
ground-aerial sensor frame pair reflects the drone's integrated physical state
over ${\approx}50$ aerial physics steps.
Ground actor states are also resolved at tick boundaries, ensuring temporal
co-registration across all sensor modalities within a single tick.
Applications requiring finer aerial state resolution should reduce the fixed
simulation delta-time accordingly; the throughput trade-off follows from
Table~\ref{tab:fps_scaling}.

\section{Performance Evaluation}
\label{sec:perf}

This section evaluates CARLA-Air under representative joint air-ground
workloads across three experiments: frame-rate and resource scaling
(Section~\ref{sec:perf:fps}), memory stability under sustained operation
(Section~\ref{sec:perf:vram}), and communication latency
(Section~\ref{sec:perf:rpc}).
Full configuration parameters and raw data are deferred to
Appendix~\ref{app:config}.

\paragraph{Reference hardware.}
All measurements are collected on a workstation equipped with an NVIDIA
RTX~A4000 (16\,GB GDDR6), AMD Ryzen~7 5800X (8-core, 4.7\,GHz), and 32\,GB
DDR4-3200, running Ubuntu~20.04~LTS.
The simulator runs in Epic quality mode with Town10HD loaded unless stated
otherwise.
All aerial experiments use the built-in SimpleFlight controller with default
PID gains.
CPU affinity and GPU power limits are left at system defaults to reflect
realistic research deployment conditions.

%% ---------------------------------------------------------------
\subsection{Benchmark Methodology}
\label{sec:perf:method}

Reliable performance measurement requires eliminating startup
transients---map-loading jitter, first-frame shader compilation, and actor
lifecycle initialization must be discarded before steady-state sampling begins.
Algorithm~\ref{alg:bench} formalizes the benchmark harness used throughout this
section.

\input{figures/alg_bench}

Each profile uses $T_w = 200$ warm-up ticks and $T_m = 2{,}000$ measurement
ticks.
VRAM is sampled every 60\,s.
All reported frame rates are the harmonic mean of $\mathbf{f}$---the
appropriate central tendency for rate quantities---with standard deviations
alongside.
The latency benchmark issues 500 warm-up calls followed by 5\,000 measurement
calls; actor spawn calls are each paired with an immediate destroy to prevent
scene-state accumulation from contaminating subsequent measurements.

%% ---------------------------------------------------------------
\subsection{Experiment~1: Frame Rate and Resource Scaling}
\label{sec:perf:fps}

\paragraph{Workload design.}
Under synchronous-mode operation, per-tick wall time is bounded by the slowest
of three concurrent contributors: the rendering thread (GPU-bound), the
ground-actor dispatch loop (CPU-bound), and the aerial physics engine
(CPU-bound, asynchronous at ${\approx}1{,}000$\,Hz on a dedicated thread).
Sensor rendering dominates at high resolution due to GPU memory bandwidth
saturation; actor population contributes via per-mesh draw calls.
These relationships motivate the workload stratification in
Table~\ref{tab:fps_scaling}.

\input{figures/tab_fps_scaling}

\paragraph{Analysis.}
The moderate joint configuration---combining ground traffic, an aerial agent,
and a full sensor suite---sustains $19.8 \pm 1.1$\,FPS, sufficient for
closed-loop policy evaluation at standard RL episode lengths.
Comparing the standalone ground baseline ($28.4$\,FPS) against the moderate
joint profile ($19.8$\,FPS) quantifies total integration overhead at
$8.6$\,FPS (30.3\%), of which $2.1$\,FPS is attributable to ground co-hosting
and the remaining $6.5$\,FPS to the aerial physics engine.
Crucially, this aerial overhead manifests entirely in CPU utilization ($54\%$
vs.\ $38\%$) rather than GPU memory: VRAM differs by only 39\,MiB between
ground-only and moderate-joint profiles.
The traffic surveillance profile retains comparable throughput ($20.1$\,FPS)
despite doubling the vehicle count, confirming that sensor rendering---not
actor population---is the dominant cost driver at $1920\times1080$.
The 3-hour endurance profile ($19.7 \pm 1.3$\,FPS) confirms that throughput
does not degrade under sustained operation.

%% ---------------------------------------------------------------
\subsection{Experiment~2: Memory Stability}
\label{sec:perf:vram}

\paragraph{Steady-state VRAM profile.}
The base map asset load accounts for approximately 3\,702\,MiB at idle
(Town10HD, Epic quality); each additional ground sensor contributes
4--10\,MiB at $1280\times720$.
A transient peak of approximately 5\,000\,MiB occurs during map loading and
resolves within 30\,s.
All steady-state profiles remain within 3\,878\,MiB, retaining approximately
12\,506\,MiB (76\%) of the 16\,GB device budget for co-located workloads such
as GPU-based policy training.

\paragraph{Leak verification.}
Table~\ref{tab:stability} summarises the 3-hour endurance run across 357 actor
spawn-and-destroy cycles.
A linear regression of VRAM against cycle index yields a slope of
$0.49$\,MiB/cycle with $R^2 = 0.11$, indicating no statistically significant
accumulation trend.
Fig.~\ref{fig:vram_timeseries} visualises the VRAM trace: the negligible
early-to-late drift ($3{,}868 \to 3{,}878$\,MiB, approximately 10\,MiB over
327 cycles) is attributable to residual render-target caching rather than
lifecycle leakage.

\input{figures/fig_vram_timeseries}
\input{figures/tab_stability}

Zero API errors and zero simulation crashes across 357 cycles validate the
robustness of the joint environment under the repeated agent-reset patterns
typical of reinforcement-learning training.

%% ---------------------------------------------------------------
\subsection{Experiment~3: Communication Latency}
\label{sec:perf:rpc}

Both simulation APIs operate within the same process on the loopback interface,
eliminating inter-process serialization overhead.
The two RPC servers use distinct wire protocols and port assignments (detailed
in Appendix~\ref{app:config}).
Table~\ref{tab:rpc_latency} reports round-trip latency for representative API
calls.

\input{figures/tab_rpc_latency}

\paragraph{Analysis.}
Lightweight state queries complete in 280--490\,$\mu$s, well below the per-tick
budget at 20\,FPS (50\,ms), confirming that API overhead does not contribute
meaningfully to frame-time variance.
Actor spawn latency (1\,850\,$\mu$s) reflects a one-time GPU synchronization
point for render-asset registration at episode reset.
Image capture latency (3\,200\,$\mu$s) covers the full sensor
pipeline---rendering, buffer readback, and serialization---and is overlapped
with the rendering thread at synchronous tick rates.
All measured values fall below the lower bound of bridge-based cross-process
synchronization costs reported in~\cite{transimhub}
(1\,000--5\,000\,$\mu$s per frame).

\paragraph{Tick-rate reconciliation.}
The aerial physics engine advances at ${\approx}1{,}000$\,Hz on its dedicated
thread, while the rendering tick runs at ${\approx}20$\,Hz under the moderate
joint workload.
Sensor callbacks read the aerial state at rendering tick boundaries, so each
ground-aerial sensor frame pair reflects the drone's integrated physical state
over ${\approx}50$ aerial physics steps.
Ground actor states are also resolved at tick boundaries, ensuring temporal
co-registration across all sensor modalities within a single tick.
Applications requiring finer aerial state resolution should reduce the fixed
simulation delta-time accordingly; the throughput trade-off follows from
Table~\ref{tab:fps_scaling}.

\section{Performance Evaluation}
\label{sec:perf}

This section evaluates CARLA-Air under representative joint air-ground
workloads across three experiments: frame-rate and resource scaling
(Section~\ref{sec:perf:fps}), memory stability under sustained operation
(Section~\ref{sec:perf:vram}), and communication latency
(Section~\ref{sec:perf:rpc}).
Full configuration parameters and raw data are deferred to
Appendix~\ref{app:config}.

\paragraph{Reference hardware.}
All measurements are collected on a workstation equipped with an NVIDIA
RTX~A4000 (16\,GB GDDR6), AMD Ryzen~7 5800X (8-core, 4.7\,GHz), and 32\,GB
DDR4-3200, running Ubuntu~20.04~LTS.
The simulator runs in Epic quality mode with Town10HD loaded unless stated
otherwise.
All aerial experiments use the built-in SimpleFlight controller with default
PID gains.
CPU affinity and GPU power limits are left at system defaults to reflect
realistic research deployment conditions.

%% ---------------------------------------------------------------
\subsection{Benchmark Methodology}
\label{sec:perf:method}

Reliable performance measurement requires eliminating startup
transients---map-loading jitter, first-frame shader compilation, and actor
lifecycle initialization must be discarded before steady-state sampling begins.
Algorithm~\ref{alg:bench} formalizes the benchmark harness used throughout this
section.

\input{figures/alg_bench}

Each profile uses $T_w = 200$ warm-up ticks and $T_m = 2{,}000$ measurement
ticks.
VRAM is sampled every 60\,s.
All reported frame rates are the harmonic mean of $\mathbf{f}$---the
appropriate central tendency for rate quantities---with standard deviations
alongside.
The latency benchmark issues 500 warm-up calls followed by 5\,000 measurement
calls; actor spawn calls are each paired with an immediate destroy to prevent
scene-state accumulation from contaminating subsequent measurements.

%% ---------------------------------------------------------------
\subsection{Experiment~1: Frame Rate and Resource Scaling}
\label{sec:perf:fps}

\paragraph{Workload design.}
Under synchronous-mode operation, per-tick wall time is bounded by the slowest
of three concurrent contributors: the rendering thread (GPU-bound), the
ground-actor dispatch loop (CPU-bound), and the aerial physics engine
(CPU-bound, asynchronous at ${\approx}1{,}000$\,Hz on a dedicated thread).
Sensor rendering dominates at high resolution due to GPU memory bandwidth
saturation; actor population contributes via per-mesh draw calls.
These relationships motivate the workload stratification in
Table~\ref{tab:fps_scaling}.

\input{figures/tab_fps_scaling}

\paragraph{Analysis.}
The moderate joint configuration---combining ground traffic, an aerial agent,
and a full sensor suite---sustains $19.8 \pm 1.1$\,FPS, sufficient for
closed-loop policy evaluation at standard RL episode lengths.
Comparing the standalone ground baseline ($28.4$\,FPS) against the moderate
joint profile ($19.8$\,FPS) quantifies total integration overhead at
$8.6$\,FPS (30.3\%), of which $2.1$\,FPS is attributable to ground co-hosting
and the remaining $6.5$\,FPS to the aerial physics engine.
Crucially, this aerial overhead manifests entirely in CPU utilization ($54\%$
vs.\ $38\%$) rather than GPU memory: VRAM differs by only 39\,MiB between
ground-only and moderate-joint profiles.
The traffic surveillance profile retains comparable throughput ($20.1$\,FPS)
despite doubling the vehicle count, confirming that sensor rendering---not
actor population---is the dominant cost driver at $1920\times1080$.
The 3-hour endurance profile ($19.7 \pm 1.3$\,FPS) confirms that throughput
does not degrade under sustained operation.

%% ---------------------------------------------------------------
\subsection{Experiment~2: Memory Stability}
\label{sec:perf:vram}

\paragraph{Steady-state VRAM profile.}
The base map asset load accounts for approximately 3\,702\,MiB at idle
(Town10HD, Epic quality); each additional ground sensor contributes
4--10\,MiB at $1280\times720$.
A transient peak of approximately 5\,000\,MiB occurs during map loading and
resolves within 30\,s.
All steady-state profiles remain within 3\,878\,MiB, retaining approximately
12\,506\,MiB (76\%) of the 16\,GB device budget for co-located workloads such
as GPU-based policy training.

\paragraph{Leak verification.}
Table~\ref{tab:stability} summarises the 3-hour endurance run across 357 actor
spawn-and-destroy cycles.
A linear regression of VRAM against cycle index yields a slope of
$0.49$\,MiB/cycle with $R^2 = 0.11$, indicating no statistically significant
accumulation trend.
Fig.~\ref{fig:vram_timeseries} visualises the VRAM trace: the negligible
early-to-late drift ($3{,}868 \to 3{,}878$\,MiB, approximately 10\,MiB over
327 cycles) is attributable to residual render-target caching rather than
lifecycle leakage.

\input{figures/fig_vram_timeseries}
\input{figures/tab_stability}

Zero API errors and zero simulation crashes across 357 cycles validate the
robustness of the joint environment under the repeated agent-reset patterns
typical of reinforcement-learning training.

%% ---------------------------------------------------------------
\subsection{Experiment~3: Communication Latency}
\label{sec:perf:rpc}

Both simulation APIs operate within the same process on the loopback interface,
eliminating inter-process serialization overhead.
The two RPC servers use distinct wire protocols and port assignments (detailed
in Appendix~\ref{app:config}).
Table~\ref{tab:rpc_latency} reports round-trip latency for representative API
calls.

\input{figures/tab_rpc_latency}

\paragraph{Analysis.}
Lightweight state queries complete in 280--490\,$\mu$s, well below the per-tick
budget at 20\,FPS (50\,ms), confirming that API overhead does not contribute
meaningfully to frame-time variance.
Actor spawn latency (1\,850\,$\mu$s) reflects a one-time GPU synchronization
point for render-asset registration at episode reset.
Image capture latency (3\,200\,$\mu$s) covers the full sensor
pipeline---rendering, buffer readback, and serialization---and is overlapped
with the rendering thread at synchronous tick rates.
All measured values fall below the lower bound of bridge-based cross-process
synchronization costs reported in~\cite{transimhub}
(1\,000--5\,000\,$\mu$s per frame).

\paragraph{Tick-rate reconciliation.}
The aerial physics engine advances at ${\approx}1{,}000$\,Hz on its dedicated
thread, while the rendering tick runs at ${\approx}20$\,Hz under the moderate
joint workload.
Sensor callbacks read the aerial state at rendering tick boundaries, so each
ground-aerial sensor frame pair reflects the drone's integrated physical state
over ${\approx}50$ aerial physics steps.
Ground actor states are also resolved at tick boundaries, ensuring temporal
co-registration across all sensor modalities within a single tick.
Applications requiring finer aerial state resolution should reduce the fixed
simulation delta-time accordingly; the throughput trade-off follows from
Table~\ref{tab:fps_scaling}.